% =====================================================================
%  Chapter 27 opener -- "Chapter at a Glance", three reading paths + the
%  ladder staircase, THE running example (the shared 8-node rooted tree R,
%  its preorder flatten, its doubled Euler tour, and a colouring), the
%  who-owns-what comparison against Ch.16 (LCA), Ch.5 (Euler/tin), Ch.17
%  (DSU on tree), Ch.22 (HLD), Ch.24 (centroid) and Ch.26 (virtual trees),
%  the recognition guide, a constraint->technique table, and notation.
%  Exports fig:c21-tree, fig:c21-taxonomy.
%  No \term{lca} / \term{tin} / \term{Euler tour} here -- they are indexed by
%  Ch.16/Ch.5; \term the NEW objects "Mo's algorithm", "block order", "toggle",
%  "offline LCA" only.
% =====================================================================

\begin{chaptertheorybox}{Chapter at a Glance}
Chapter~5 (\chref{c24}) taught you to \emph{flatten} a tree: one DFS fixes
\(\tin,\tout\), and the preorder array turns every subtree into a contiguous
range.  Chapter~16 (\chref{c11}) taught you to \emph{climb}: binary lifting
answers \(\lca\) and \(\dist\) online in \(O(\log n)\) (\(\thmref{c11-distf}\)).
Both are single-query tools -- order the queries however you like, each is
answered on its own.

This chapter is about the opposite regime: a \emph{huge batch of offline
queries} over the \emph{same} tree, where the statistic you maintain has
\textbf{no inverse}.  "How many distinct colours are on this path?" you cannot
answer by subtracting two prefixes -- there is no \(\mathrm{pref}(r)-
\mathrm{pref}(l-1)\) for "distinct".  When add/remove of one element is cheap
but there is no undo-by-arithmetic, you stop trying to be online and instead
\emph{reorder} the queries so that consecutive ones differ by a single element:
that is Mo's algorithm, a \(\bigO{(n+q)\sqrt n}\) sweep with a
\term{block order}.  And when the queries only ask for \(\lca\), Tarjan's
offline DFS with a disjoint-set union answers all \(q\) of them in
\(\bigO{(n+q)\alpha(n)}\) -- no lifting table at all.

\textbf{You will learn to}
\begin{itemize}[nosep]
  \item run Mo's \term{block order}: sort \(q\) ranges by \((\lfloor L/B\rfloor, R)\),
        move two pointers with \(O(1)\) add/remove, and bound the work by
        \(\bigO{(n+q)\sqrt n}\) with \(B\approx\sqrt n\)
        (\(\secref{c21-moengine}\), solved CF 617E \emph{XOR and Favorite Number}); \hfill\lv{3}--\lv{4}
  \item answer \term{Mo on subtrees}: flatten subtree\((v)\) to \([\tin(v),\tout(v)]\)
        and reuse the engine; carry root-path statistics on the DFS stack
        (\(\secref{c21-subtree}\), solved CF 463E \emph{Caisa and Tree}); \hfill\lv{4}
  \item answer \term{Mo on paths}: keep the doubled Euler tour, \term{toggle}
        membership on each occurrence, and split at \(\lca\) with the LCA-parity
        special case (\(\secref{c21-path}\)); \hfill\lv{5}
  \item compute \term{offline LCA} with Tarjan's DFS+DSU in \(\bigO{(n+q)\alpha}\),
        and know when it beats Ch.~16 lifting (\(\secref{c21-tarjan}\)); \hfill\lv{4}--\lv{5}
  \item read a statement and pick Mo over Fenwick / HLD / DSU-on-tree / virtual
        tree / centroid exactly when the statistic has no inverse and the queries
        are offline (\(\secref{c21-boundary}\)). \hfill\lv{5}
\end{itemize}

\textbf{Prerequisites.} Chapter~5 (\chref{c24}) -- \(\tin/\tout\) and the
flatten-to-array trick: the range a subtree becomes.  Chapter~16 (\chref{c11}) --
\(\lca\), used to split a path (27.3) and the online target Tarjan replaces (27.4).
Chapter~13 (\chref{c10}) -- the intuition that a \(\bigO{n}\) fold per query is too
slow when there are \(q\) of them.  Chapter~17 (\chref{c28}) -- DSU on tree, the
\emph{other} offline answer to subtree statistics (27.5 contrasts them).
Chapter~26 (\chref{c30}) -- the \textbf{previous chapter}; its method boundary hands
the "many overlapping offline node-set queries, non-invertible statistic" case here.
\end{chaptertheorybox}

% ---------------- Reading paths + ladder position ----------------
\begin{chaptertheorybox}{How to Read This Chapter (three paths)}
Level badges \(\lv3\)--\lv{5}: \(\lv3\) is running the block order by hand on a
toy array and counting pointer moves; \(\lv4\) is the two coded solves --
array Mo (617E) and the root-path DFS-stack sweep (463E); \(\lv{5}\) is the doubled-tour
\emph{toggle} for paths and Tarjan's DFS+DSU, capped by the Summit (CF 1585E).

\textbf{Engine path (\(\lv3\)--\(\lv4\), ``make the sweep'').}
\secref{c21-moengine} in full -- block size, the two-pointer move, the
\(\bigO{(n+q)\sqrt n}\) bound, solved CF 617E -- then \secref{c21-subtree}, where
the same engine runs on the flattened subtree ranges (solved CF 463E).

\textbf{Tree-power path (\(\lv4\)--\(\lv{5}\), ``a statistic with no inverse, on
a path'').} \secref{c21-subtree} for the flatten, \secref{c21-path} for the
doubled Euler tour and the toggle + LCA split, then the Summit \(\probref{c21-summit}\),
where a frequency-of-frequency bucket answers offline root-path queries.

\textbf{Boundary path (\(\lv4\)--\(\lv{5}\), ``is this even a Mo job?'').}
\secref{c21-tarjan} (offline LCA vs Ch.~16 lifting) and \secref{c21-boundary}:
the tell (offline, overlapping, no inverse) and the look-alikes that are not --
invertible ranges for Fenwick, online paths for HLD, per-set folds for a virtual
tree, small marked sets for DSU-on-tree, all-pairs for centroid.
\end{chaptertheorybox}

\begin{figure}[H]
\centering
\begin{tikzpicture}[diagram, x=8.5mm, y=6.5mm]
  \foreach \st/\lab in {1/lv1,2/lv2,3/lv3,4/lv4,5/lv5,6/lv6} {
    \fill[\lab!22, draw=\lab, line width=0.8pt] (\st-1,0) rectangle (\st,\st*0.52);
    \node[font=\scriptsize\bfseries, text=\lab] at (\st-0.5,\st*0.52-0.42) {L\st};}
  \node[annot, anchor=south west, font=\scriptsize]
    at (0.05,0.62) {prereq (Ch.5/16): tin/tout flatten, lca -- the range and the split};
  \node[annot, anchor=south west, font=\scriptsize]
    at (2.05,1.66) {27.1 block order, two pointers, the (n+q)sqrt n bound};
  \node[annot, anchor=south west, font=\scriptsize]
    at (3.05,2.18) {27.1 solved CF 617E (array Mo) | 27.2 solved CF 463E (root path)};
  \node[annot, anchor=south west, font=\scriptsize, text=lv5]
    at (4.05,2.70) {27.3 doubled tour + toggle | 27.4 Tarjan offline LCA};
  \node[annot, anchor=south west, font=\scriptsize, text=lv6]
    at (5.05,3.22) {27.5 the pick; Summit CF 1585E (offline root-path frequency)};
\end{tikzpicture}
\caption{Chapter~27 on the six-step ladder: hand-running the block order and
counting moves is \(\lv3\); the coded array-Mo and DFS-stack-root-path solves (617E,
463E) are \(\lv4\); the doubled-tour toggle for paths and Tarjan's offline LCA
are \(\lv{5}\); the Summit (an offline root-path frequency query) is step~6.}
\end{figure}

% ---------------- The running example ----------------
\subsection*{The Running Example: One Tree, Three Flattenings}

Every section in this chapter runs on the same rooted tree \(R\) on \(n=8\)
nodes, fixed once here and authoritative (if a later trace disagrees, this
figure wins).  Rooted at \(1\); parent-to-child edges, 1-based:
\[
  1\!\to\!2,\quad 1\!\to\!3,\quad 2\!\to\!4,\quad 2\!\to\!5,\quad
  3\!\to\!6,\quad 3\!\to\!7,\quad 4\!\to\!8 .
\]
A left-to-right DFS gives the \term{preorder flatten} \(P\) (entry times
\(\tin\)): \(P=[\,1,2,4,8,5,3,6,7\,]\), so \(\tin(1){=}1\), \(\tin(2){=}2\),
\(\tin(4){=}3\), \(\tin(8){=}4\), \(\tin(5){=}5\), \(\tin(3){=}6\),
\(\tin(6){=}7\) and \(\tin(7){=}8\); subtree\((2)=\{2,4,8,5\}\) is therefore the
single contiguous range \(P[2..5]\) (\(\secref{c21-subtree}\)).  Depths are
\(\depth(1){=}0\), \(\depth(2){=}\depth(3){=}1\),
\(\depth(4){=}\depth(5){=}\depth(6){=}\depth(7){=}2\) and \(\depth(8){=}3\).

For path queries we instead keep the \emph{doubled} Euler tour \(E\) (a node is
written when we enter \emph{and} when we return to it), each node appearing twice:
\[
  E=[\,1,2,4,8,8,4,5,5,2,3,6,6,7,7,3,1\,] ,
\]
with \(st(v)\) the first index of \(v\) and \(en(v)\) its second, listed for the
nodes in preorder order \((1,2,4,8,5,3,6,7)\):
\(st{=}(1,2,3,4,7,10,11,13)\) and \(en{=}(16,9,6,5,8,15,12,14)\)
(\(\secref{c21-path}\)).  Finally a \emph{colouring}
\(\mathit{val}(v)=2,1,3,1,2,3,1,2\) for \(v=1..8\) makes "how many distinct colours on
this range/path?" a genuinely non-invertible statistic -- the reason Mo, and not
a Fenwick prefix difference, is the right tool.

\begin{figure}[H]
\centering
\begin{tikzpicture}[diagram, scale=0.95, transform shape,
   tnode/.style={circle, draw=primaryblue!70, fill=primaryblue!8, line width=0.7pt,
                 minimum size=6.6mm, font=\scriptsize\bfseries, inner sep=1pt},
   sub/.style={circle, draw=difforange, fill=difforange!16, line width=1.0pt,
             minimum size=6.6mm, font=\scriptsize\bfseries, inner sep=1pt},
   val/.style={font=\tiny, text=warningred}]
  \node[tnode] (n1) at (0,0) {1};
  \node[sub]   (n2) at (-2.4,-1.5) {2};
  \node[tnode] (n3) at (2.4,-1.5) {3};
  \node[sub]   (n4) at (-3.6,-3.0) {4};
  \node[sub]   (n5) at (-1.4,-3.0) {5};
  \node[tnode] (n6) at (1.4,-3.0) {6};
  \node[tnode] (n7) at (3.6,-3.0) {7};
  \node[sub]   (n8) at (-3.6,-4.5) {8};
  \draw (n1)--(n2) (n1)--(n3) (n2)--(n4) (n2)--(n5) (n3)--(n6) (n3)--(n7) (n4)--(n8);
  \node[val] at (0.45,0.35) {\tt 2};   \node[val] at (-1.95,-1.15) {\tt 1};
  \node[val] at (2.85,-1.15) {\tt 3};  \node[val] at (-4.05,-2.65) {\tt 1};
  \node[val] at (-0.95,-2.65) {\tt 2}; \node[val] at (1.85,-2.65) {\tt 3};
  \node[val] at (4.05,-2.65) {\tt 1};  \node[val] at (-4.05,-4.15) {\tt 2};
  \node[annot, font=\scriptsize, anchor=north west] at (-4.9,0.9)
    {\tt val(v) in red; orange = subtree(2) \{2,4,8,5\}=P[2..5]};
\end{tikzpicture}
\caption{The running example \(R\) (8 nodes, root 1).  Orange nodes are the
subtree of \(2\), which the preorder flatten makes the single contiguous range
\(P[\tin(2){=}2\,..\,\tout(2){=}5]\) (\(\secref{c21-subtree}\)).  Red labels are
\(\mathit{val}(v)\), the colour whose ``distinct count on a range/path'' statistic has no
inverse -- the whole reason this chapter exists.}
\label{fig:c21-tree}
\end{figure}

% ---------------- Which job, cost chart ----------------
\subsection*{Which Job? (the question-shape chart)}

\begin{center}
\small
\renewcommand{\arraystretch}{1.25}
\setlength{\tabcolsep}{3pt}
\begin{tabular}{@{}p{0.31\linewidth} p{0.17\linewidth} p{0.28\linewidth} >{\raggedright\arraybackslash}p{0.14\linewidth}@{}}
\toprule
\textbf{Question} & \textbf{Technique} & \textbf{Why not a Fenwick/segment prefix} & \textbf{Where} \\
\midrule
count subarrays with \(\mathit{xor}=k\), offline & array Mo, toggle prefix-xor freq & ``equal to \(k\) among freq of xors'' has no per-index inverse & \secref{c21-moengine} \\
distinct colours in a subtree, offline & flatten + Mo & \(\sum\!\mathrm{distinct}\) is not pref\((r)-\)pref\((l{-}1)\) & \secref{c21-subtree} \\
add to a node, sign-flip down its subtree & preorder + parity BIT & this one \emph{is} invertible: it goes to Fenwick & \secref{c21-subtree} \\
distinct colours on a path \(u\)-\(v\), offline & doubled tour + toggle & a path is two root-paths minus overlap; ranges do not reach it & \secref{c21-path} \\
answer \(q\) LCAs offline, no lifting table & Tarjan DFS+DSU & lifting is per-query \(O(\log n)\); Tarjan is \(O(\alpha)\) batch & \secref{c21-tarjan} \\
an \emph{online} path range with updates & \emph{go to HLD / LCT} & Mo is offline only; a moving/queried-now model breaks it & \chref{c16} \\
\bottomrule
\end{tabular}
\end{center}

\begin{insightbox}[One promise decides everything: reorder so neighbours differ by one]
A statistic with no inverse (distinct count, \(\sum \mathrm{freq}^2\), "how many
values appear \(\ge l\) times") cannot be answered from two prefixes, but
\emph{add/remove of one element} is cheap.  Mo's whole idea is to exploit that
asymmetry: instead of \(q\) independent answers, walk a window across the array
visiting every query range, ordering them so that the window moves as little as
possible.  Block ordering bounds the total movement to \(\bigO{(n+q)\sqrt n}\).
Every tree version is the same sweep on a smarter array: the preorder flatten
(subtrees), the doubled Euler tour with a \emph{toggle} (paths).  The one
chapter idea: \emph{when you can add cheaply but cannot subtract, change the
\emph{order} of the questions so you only ever add or drop one node.}
\end{insightbox}

% ---------------- Taxonomy chart ----------------
\begin{figure}[H]
\centering
\begin{tikzpicture}[diagram, scale=0.95, transform shape,
    mnode/.style={circle, draw=#1, line width=0.9pt, fill=#1!8, minimum size=9mm,
                  font=\scriptsize\bfseries, text=#1!60!black}]
  \node[mnode=primaryblue] (a) at (0,0)   {order};
  \node[annot, font=\scriptsize, anchor=west, text width=9.8cm, align=left] at (1.7,0) {block-sort the q offline ranges by (left block, right); a sweep answers all of them. B approx sqrt n, total O((n+q) sqrt n) (27.1)};
  \node[mnode=successgreen] (b) at (0,-1.1) {flatten};
  \node[annot, font=\scriptsize, anchor=west, text width=9.8cm, align=left] at (1.7,-1.1) {a subtree(v) is the contiguous preorder range [tin, tout]; run the same block sweep on it -- and when the queried sets are root paths, the DFS stack itself is the window (27.2)};
  \node[mnode=difforange]   (c) at (0,-2.2) {toggle};
  \node[annot, font=\scriptsize, anchor=west, text width=9.8cm, align=left] at (1.7,-2.2) {a path is not contiguous, so use the doubled Euler tour E[1..2n]; a node is in the path iff it appears exactly once in the current window; split at lca and add it by hand when neither endpoint is it (27.3)};
  \node[mnode=softpurple]   (d) at (0,-3.3) {offline};
  \node[annot, font=\scriptsize, anchor=west, text width=9.8cm, align=left] at (1.7,-3.3) {for LCA-only queries, Tarjan's single DFS with DSU answers all q in O((n+q) alpha) -- a WHITE/GREY/BLACK colouring, ancestor[find(u)] at the moment the second endpoint turns BLACK (27.4)};
  \node[mnode=warningred]   (e) at (0,-4.4) {reject};
  \node[annot, font=\scriptsize, anchor=west, text width=9.9cm, align=left] at (1.7,-4.4) {invertible range -> Fenwick (Ch.5/10); online path -> HLD (Ch.22) / LCT (Ch.25); per-set fold -> virtual tree (Ch.26); small marked set -> DSU-on-tree (Ch.17); all-pairs distance -> centroid (Ch.24) (27.5)};
\end{tikzpicture}
\caption{The chapter's five-state taxonomy.  \emph{order} is the array Mo sweep;
\emph{flatten} lifts it to subtrees; \emph{toggle} lifts it to paths via the
doubled tour; \emph{offline} is Tarjan's batched LCA; \emph{reject} hands the job
to the right neighbour.  Each solved problem sits on one row: CF 617E on
\emph{order}, CF 463E riding the recursion stack beside \emph{flatten}, the
Summit CF 1585E on \emph{toggle}+a frequency bucket.}
\label{fig:c21-taxonomy}
\end{figure}

% ---------------- Who owns what ----------------
\subsection*{Who Owns What (so nothing is learned twice)}

\begin{center}
\small
\renewcommand{\arraystretch}{1.25}
\setlength{\tabcolsep}{3pt}
\begin{tabular}{@{}p{0.28\linewidth} p{0.24\linewidth} p{0.44\linewidth}@{}}
\toprule
\textbf{Idea} & \textbf{Owner} & \textbf{This chapter's share} \\
\midrule
\(\tin/\tout\), subtree-as-range flatten & Ch.~5 (\chref{c24}) & the range Mo sweeps over; the single-visit flatten is reused, not re-taught \\
\(\lca\), \(\dist\), binary lifting & Ch.~16 (\chref{c11}) & 27.3 calls \(\lca\) to split a path; 27.4 offers Tarjan as an offline \emph{alternative}, never re-deriving lifting \\
invertible range add / point query & Ch.~5 / Ch.~10 Fenwick & 27.2 points back to it (Ch.~18 solves CF 383C with it, \chref{c14}); 27.5 routes every invertible query back to it \\
small-to-large / DSU on tree & Ch.~17 (\chref{c28}) & the rival offline answer for subtree stats (owns CF 375D); 27.5 draws the Mo-vs-DSU line \\
online path / subtree ranges & Ch.~22 HLD (\chref{c16}) & explicitly out; Mo is offline, a live path query with updates is HLD's job \\
per-marked-set folds on a fixed tree & Ch.~26 virtual tree (\chref{c30}) & the previous chapter hands \emph{overlapping} offline node-set queries to Mo \\
\bottomrule
\end{tabular}
\end{center}

% ---------------- Recognition guide ----------------
\begin{recognitionbox}[Pattern card -- ``a big offline batch of subtree/path queries, no inverse'']
\reckey{statements that fix one tree and hand you a \emph{large batch} of queries,
each asking a statistic over a subtree or a path, where the statistic has
\emph{no inverse}: number of distinct colours, \(\sum\mathrm{freq}^2\), ``how many
values occur \(\ge l\) times'', ``the \(k\)-th most frequent''.  The tell is
\(\sum\) over the ranges does not telescope, but adding or removing a single node
changes the answer in \(O(1)\).  All queries are known up front (offline) and no
query depends on the answer to an earlier one.}
\recuse{
  \item distinct-colour / frequency statistics over subtree ranges -- flatten + Mo (\(\secref{c21-subtree}\));
  \item the same over paths -- doubled Euler tour + toggle (\(\secref{c21-path}\));
  \item a pure batch of \(q\) LCA queries with no updates -- Tarjan offline (\(\secref{c21-tarjan}\)).
}
\recnotuse{
  \item the statistic \emph{does} have an inverse (a sum, an xor, a count of one fixed value) -- Euler + Fenwick / prefix array (\chref{c24});
  \item queries are \emph{online} or interleaved with updates -- HLD (\chref{c16}) or link-cut (\chref{c23});
  \item each query names a tiny marked set to fold over -- virtual tree (\chref{c30});
  \item you recompute a per-node subtree value for the \emph{whole} tree once -- tree DP / DSU-on-tree (\chref{c10}, \chref{c28});
  \item the answer is a distance distribution over all pairs -- centroid (\chref{c18}).
}
One-line test: \emph{offline batch + add/remove is cheap + the statistic does not
telescope -- then Mo; if it telescopes, use Fenwick, if it is online, use HLD.}
\end{recognitionbox}

% ---------------- Constraint -> technique ----------------
\subsection*{Reading the Constraints (which technique the limits want)}

\begin{center}
\small
\renewcommand{\arraystretch}{1.25}
\setlength{\tabcolsep}{3pt}
\begin{tabular}{@{}p{0.34\linewidth} p{0.40\linewidth} p{0.20\linewidth}@{}}
\toprule
\textbf{Constraint signature} & \textbf{What it rules in/out} & \textbf{Technique} \\
\midrule
\(n,q\le10^5\), queries known up front, ask "distinct / \(k\)-th / \(\ge l\)" over a range & telescoping prefix fails; a per-query fold is \(\bigO{nq}\) & \secref{c21-moengine} \\
subtree of \(v\), offline, non-invertible statistic & the subtree is a contiguous range once flattened & \secref{c21-subtree} \\
path \(u\)-\(v\), offline, non-invertible statistic & not contiguous -- doubled tour + toggle & \secref{c21-path} \\
\(q\) LCA queries, offline, no updates, want linear-ish total & no lifting table needed; one DFS+DSU & \secref{c21-tarjan} \\
updates interleaved, or queries must be answered online & Mo needs to see all queries first & \chref{c16} (HLD) \\
\bottomrule
\end{tabular}
\end{center}

% ---------------- Notation ----------------
\subsection*{Notation Used Only Here}

The fixed tree is \(R=(V,E)\), \(|V|=n\), rooted at \(1\).  \(P[\,]\) is the
preorder flatten (\(P[\tin(v)]=v\)); subtree\((v)\) is the range
\([\tin(v),\tout(v)]\).  \(E[\,]\) is the \emph{doubled} Euler tour of length
\(2n\), with \(st(v)\) and \(en(v)\) the two positions where \(v\) appears
(\(st\) on entry, \(en\) on exit).  A Mo window is the index interval
\([L,R]\) of the current array (either \(P\) or \(E\)); \texttt{add}/\texttt{rm}
are the \(O(1)\) steps that insert or delete the element at one end.  \(\mathit{val}(v)\)
is the colour/value at node \(v\); \(\mathit{cnt}[c]\) is how many times colour
\(c\) currently appears in the window; a \(\mathit{resp}\) structure (e.g.\
\(\mathit{bucket}[f]\)) holds the derived statistic so \texttt{add}/\texttt{rm}
stay \(O(1)\).  \(B\) is the Mo block size.  \(\alpha(\cdot)\) is the inverse
Ackermann function (Tarjan's DSU).  \(\depth,\tin,\tout,\lca,\dist\) are
Ch.~5/Ch.~16's, used unchanged.  Every section runs on the \(R\), \(P\), \(E\),
and \(\mathit{val}\) above; if a table here and a trace there disagree, this opener is
authoritative.
